\chapter{\texorpdfstring{Derived Relative Mealy Program Objects and Contextual Dynamic Proarrow Logic}{Derived Relative Mealy Program Objects and Contextual Dynamic Proarrow Logic}}
\label{ch:derived-relative-mealy-program-objects}

\section{Orientation: contextualizing derived Mealy morphisms}
\label{sec:derived-relative-orientation}

In Chapter~\ref{ch:derived-mealy-morphisms}, a derived Mealy morphism from
\(\mathbb A\) to \(\mathbb B\) was defined as a derived correspondence
\[
\chi:\mathbb X\to \mathbb A\times \mathbb B
\]
in the localized homotopy theory of internal categories. We write its two
components as
\[
\pi:\mathbb X\to\mathbb A,
\qquad
\Omega:\mathbb X\to\mathbb B.
\]
The object \(\mathbb X\) is the derived execution object, \(\pi\) records
input semantics, and \(\Omega\) is the derived \(\mathbb B\)-valued output
\(1\)-cocycle.

The purpose of this chapter is to develop the relative, contextual version of
this construction. The strict relative theory begins with a Mealy morphism
\[
\Phi=(P,\delta,\omega):\mathbb A\rightsquigarrow\mathbb B
\]
and a downstream right \(\mathbb B\)-action
\[
Y\to B_0.
\]
It forms the coupled state object
\[
P\times_{B_0}Y
\]
and the relative program category
\[
\mathsf{Prog}_{\Phi,Y}
=
\int_{\mathbb A}\Phi^\ast Y.
\]

In the derived theory, the strict pullback
\[
P\times_{B_0}Y
\]
is replaced by the homotopy-correct pullback of the derived output cocycle
against a downstream behaviour
\[
\rho:\mathbb Y\to\mathbb B.
\]
Thus the fundamental contextual execution object is
\[
\mathbb X\times_{\mathbb B}\mathbb Y.
\]

The guiding principle is:
\[
\boxed{
\text{derived contextual Mealy logic is dynamic proarrow logic on the
homotopy pullback of an output cocycle with a downstream behaviour.}
}
\]

There is a minor orientation issue which must be fixed explicitly. In the
derived Mealy chapter, a strict Mealy morphism determines the action category
\[
P\rtimes\mathbb A
\]
whose arrows are written
\[
\delta(a,x)\longrightarrow x.
\]
This is the contravariant, or backward, execution orientation. By contrast,
dynamic program categories are often displayed in the forward execution
orientation
\[
x\longrightarrow \delta(a,x).
\]
In this chapter the intrinsic derived object will be the backward execution
object
\[
\mathsf{Exec}^h_{\chi,\rho}.
\]
Whenever a forward dynamic-program orientation is desired, it is obtained by
passing to an opposite representative. Thus, when opposites are available, one
may write
\[
\mathsf{Prog}^h_{\chi,\rho}
:=
\left(\mathsf{Exec}^h_{\chi,\rho}\right)^{\mathrm{op}}.
\]
The derived construction itself is made before choosing this forward
orientation.

\section{Homotopical setting and conventions}
\label{sec:derived-relative-homotopical-setting}

Throughout this chapter, let
\[
\mathcal E
\]
be a Grothendieck topos. Let
\[
\mathcal M:=\mathbf{Cat}(\mathcal E)_{JT}
\]
be the Joyal--Tierney model category of internal categories in \(\mathcal E\),
and let
\[
\mathcal C
:=
\mathcal M_\infty
:=
L_{DK}(\mathcal M,W_{JT})
\]
be its Dwyer--Kan localization.

All products, pullbacks, slices, mapping spaces, and correspondences in this
chapter are taken in the \(\infty\)-category
\[
\mathcal C,
\]
unless explicitly stated otherwise. Point-set internal categories and internal
functors are used only as representatives of objects and morphisms in
\(\mathcal C\).

Recall from Chapter~\ref{ch:derived-mealy-morphisms} that the derived Mealy
moduli theory from \(\mathbb A\) to \(\mathbb B\) is
\[
\mathsf{Mealy}^h(\mathbb A,\mathbb B)
:=
\mathcal C/(\mathbb A\times\mathbb B).
\]
An object
\[
\chi:\mathbb X\to\mathbb A\times\mathbb B
\]
is a derived Mealy morphism
\[
\mathbb A\rightsquigarrow^h\mathbb B.
\]
By the product universal property, it is equivalently a pair
\[
\pi:\mathbb X\to\mathbb A,
\qquad
\Omega:\mathbb X\to\mathbb B.
\]

\begin{remark}
\label{rem:derived-relative-not-localization-only}
The derived construction is not merely the localization of the external
category of strict Mealy machines. A strict Mealy morphism first determines a
point-set fibrational representative
\[
\mathbb A
\xleftarrow{\pi_P}
P\rtimes\mathbb A
\xrightarrow{\Omega_P}
\mathbb B.
\]
The derived Mealy morphism represented by it is the induced object
\[
(P\rtimes\mathbb A\to\mathbb A\times\mathbb B)
\]
of
\[
\mathcal C/(\mathbb A\times\mathbb B).
\]
Thus the strict state object \(P\) is not itself the derived execution object;
the derived execution object is represented by the action category
\(P\rtimes\mathbb A\).
\end{remark}

\begin{definition}[Correspondence notation]
\label{def:derived-relative-corr-notation}
For objects
\[
\mathbb X,\mathbb Y\in\mathcal C,
\]
we write
\[
\operatorname{Corr}(\mathcal C)(\mathbb X,\mathbb Y)
:=
\mathcal C/(\mathbb X\times\mathbb Y)
\]
for the \(\infty\)-category of correspondences from \(\mathbb X\) to
\(\mathbb Y\). Thus a correspondence
\[
\mathbb X\nrightarrow\mathbb Y
\]
is a span
\[
\mathbb X
\longleftarrow
\mathbb R
\longrightarrow
\mathbb Y
\]
in \(\mathcal C\). Composition is by pullback in \(\mathcal C\).
\end{definition}

\section{Derived downstream behaviours}
\label{sec:derived-downstream-behaviours}

\begin{definition}[Derived downstream behaviour]
\label{def:derived-downstream-behaviour}
Let
\[
\mathbb B\in\mathcal C.
\]
The \(\infty\)-category of \textbf{derived downstream \(\mathbb B\)-behaviours}
is the slice
\[
\mathsf{Beh}^h(\mathbb B)
:=
\mathcal C/\mathbb B.
\]
An object
\[
\rho:\mathbb Y\to\mathbb B
\]
of this slice is called a derived downstream \(\mathbb B\)-behaviour.
\end{definition}

A derived downstream behaviour should be read as a homotopy-coherent system
whose observable behaviour is classified by \(\mathbb B\). It may represent
ordinary states, executions, refinements of executions, equivalences between
executions, and higher coherences.

Strict right actions give a supply of point-set representatives.

\begin{proposition}[Strict actions as derived downstream behaviours]
\label{prop:strict-actions-derived-behaviours}
Let
\[
(Y\xrightarrow{r}B_0,\sigma)
\]
be a strict right action of an internal category
\[
\mathbb B.
\]
Then its backward action category
\[
Y\rtimes\mathbb B
\]
has object object \(Y\), arrow object
\[
B_1\times_{t_B,B_0,r}Y,
\]
and arrows
\[
\sigma(b,y)\longrightarrow y.
\]
The projection
\[
\pi_Y:Y\rtimes\mathbb B\to\mathbb B
\]
is a split internal discrete fibration, hence a Joyal--Tierney fibration. It
therefore represents an object
\[
\pi_Y:Y\rtimes\mathbb B\to\mathbb B
\]
of
\[
\mathcal C/\mathbb B.
\]
\end{proposition}

\begin{proof}
This is the action-category construction of
Chapter~\ref{ch:derived-mealy-morphisms} applied to the right action
\[
\sigma:B_1\times_{t_B,B_0,r}Y\to Y.
\]
The arrow \((b,y)\) has source \(\sigma(b,y)\) and target \(y\), and the
projection to \(\mathbb B\) sends it to \(b\). The defining pullback square
for the arrow object is precisely the discrete-fibration square. The strict
action unit and associativity laws give the split compatibility with
identities and composition. By the result that split internal discrete
fibrations are Joyal--Tierney fibrations, \(\pi_Y\) is a fibrational point-set
representative of an object of the derived slice.
\end{proof}

\begin{remark}
\label{rem:derived-behaviours-broader-than-actions}
The slice
\[
\mathcal C/\mathbb B
\]
is broader than the strict category of right \(\mathbb B\)-actions. Strict
right actions give rigid fibrational representatives, but a general derived
downstream behaviour is an arbitrary object over \(\mathbb B\) in the
localized homotopy theory.
\end{remark}

\section{Contextual restriction along a derived Mealy morphism}
\label{sec:contextual-restriction-derived-mealy}

Let
\[
\chi:\mathbb X\to\mathbb A\times\mathbb B
\]
be a derived Mealy morphism, with components
\[
\pi:\mathbb X\to\mathbb A,
\qquad
\Omega:\mathbb X\to\mathbb B.
\]

\begin{definition}[Contextual restriction]
\label{def:contextual-restriction-derived}
The \textbf{contextual restriction functor}
\[
\chi^\ast:
\mathcal C/\mathbb B
\longrightarrow
\mathcal C/\mathbb A
\]
is defined on a downstream behaviour
\[
\rho:\mathbb Y\to\mathbb B
\]
by
\[
\chi^\ast(\rho)
:=
\left(
\mathbb X\times_{\mathbb B}\mathbb Y
\xrightarrow{\operatorname{pr}_{\mathbb X}}
\mathbb X
\xrightarrow{\pi}
\mathbb A
\right),
\]
where the pullback is taken over
\[
\Omega:\mathbb X\to\mathbb B
\qquad
\text{and}
\qquad
\rho:\mathbb Y\to\mathbb B.
\]
Equivalently,
\[
\chi^\ast
=
\pi_!\Omega^\ast,
\]
where
\[
\Omega^\ast:\mathcal C/\mathbb B\to\mathcal C/\mathbb X
\]
is pullback and
\[
\pi_!:\mathcal C/\mathbb X\to\mathcal C/\mathbb A
\]
is postcomposition with \(\pi\).
\end{definition}

\[
\begin{tikzcd}[column sep=large, row sep=large]
\mathbb X\times_{\mathbb B}\mathbb Y
\arrow[r, "\operatorname{pr}_{\mathbb Y}"]
\arrow[d, "\operatorname{pr}_{\mathbb X}"']
&
\mathbb Y
\arrow[d, "\rho"]
\\
\mathbb X
\arrow[r, "\Omega"']
\arrow[dr, "\pi"']
&
\mathbb B
\\
&
\mathbb A
\end{tikzcd}
\]

\begin{theorem}[Derived contextual restriction]
\label{thm:derived-contextual-restriction}
Every derived Mealy morphism
\[
\chi:\mathbb A\rightsquigarrow^h\mathbb B
\]
induces a functor
\[
\chi^\ast:
\mathsf{Beh}^h(\mathbb B)\to\mathsf{Beh}^h(\mathbb A).
\]
This construction is invariant under equivalence of derived Mealy
representatives.
\end{theorem}

\begin{proof}
The assignment
\[
\rho:\mathbb Y\to\mathbb B
\longmapsto
\mathbb X\times_{\mathbb B}\mathbb Y\to\mathbb A
\]
is the composite of pullback along \(\Omega\) and postcomposition along
\(\pi\). Both operations are functorial in an \(\infty\)-category with
pullbacks.

If \(\chi\) is replaced by an equivalent object of
\[
\mathcal C/(\mathbb A\times\mathbb B),
\]
then the corresponding diagrams over \(\mathbb B\) are equivalent. Pullbacks
in \(\mathcal C\) are invariant under equivalence of diagrams, so the induced
contextual restrictions are equivalent.
\end{proof}

\begin{remark}[Operational reading]
\label{rem:contextual-restriction-operational-reading}
The functor
\[
\chi^\ast
\]
places the derived Mealy transducer in front of the downstream behaviour
\[
\rho:\mathbb Y\to\mathbb B.
\]
An execution of the resulting contextual system is a compatible pair:
an execution of the Mealy transducer and an execution of the downstream
behaviour whose intermediate \(\mathbb B\)-behaviours agree coherently.
\end{remark}

\section{Derived relative execution objects}
\label{sec:derived-relative-execution-objects}

\begin{definition}[Derived relative execution object]
\label{def:derived-relative-execution-object}
Let
\[
\chi=(\pi,\Omega):\mathbb X\to\mathbb A\times\mathbb B
\]
be a derived Mealy morphism and let
\[
\rho:\mathbb Y\to\mathbb B
\]
be a derived downstream behaviour. The \textbf{derived relative execution
object} of \(\chi\) in the context \(\rho\) is
\[
\mathsf{Exec}^h_{\chi,\rho}
:=
\mathbb X\times_{\mathbb B}\mathbb Y.
\]
\end{definition}

It comes equipped with canonical maps
\[
\operatorname{pr}_{\mathbb X}:
\mathsf{Exec}^h_{\chi,\rho}\to\mathbb X,
\qquad
\operatorname{pr}_{\mathbb Y}:
\mathsf{Exec}^h_{\chi,\rho}\to\mathbb Y.
\]

\begin{definition}[Input and downstream projections]
\label{def:derived-input-downstream-projections}
The \textbf{derived input projection} is
\[
I^h_{\chi,\rho}:
\mathsf{Exec}^h_{\chi,\rho}
\to
\mathbb A
\]
defined by
\[
I^h_{\chi,\rho}
=
\pi\operatorname{pr}_{\mathbb X}.
\]
The \textbf{derived downstream comparison} is
\[
\Lambda^h_{\chi,\rho}:
\mathsf{Exec}^h_{\chi,\rho}
\to
\mathbb Y
\]
defined by
\[
\Lambda^h_{\chi,\rho}
=
\operatorname{pr}_{\mathbb Y}.
\]
\end{definition}

\[
\begin{tikzcd}[column sep=large]
&
\mathsf{Exec}^h_{\chi,\rho}
\arrow[dl, "I^h_{\chi,\rho}"']
\arrow[dr, "\Lambda^h_{\chi,\rho}"]
&
\\
\mathbb A
&
&
\mathbb Y
\arrow[d, "\rho"]
\\
&
&
\mathbb B
\end{tikzcd}
\]

\begin{remark}[Forward program orientation]
\label{rem:forward-program-orientation-derived}
The object
\[
\mathsf{Exec}^h_{\chi,\rho}
\]
is the invariant derived execution object. In point-set representatives coming
from strict Mealy machines, its arrows are naturally oriented
target-to-source with respect to input arrows. If a forward dynamic-program
orientation is required, one passes to an opposite representative and writes
formally
\[
\mathsf{Prog}^h_{\chi,\rho}
:=
\left(\mathsf{Exec}^h_{\chi,\rho}\right)^{\mathrm{op}},
\]
when the opposite operation is available in the relevant model or localized
setting.
\end{remark}

\begin{proposition}[Labelled contextual execution object]
\label{prop:labelled-contextual-execution-object}
For every derived Mealy morphism
\[
\chi:\mathbb A\rightsquigarrow^h\mathbb B
\]
and downstream behaviour
\[
\rho:\mathbb Y\to\mathbb B,
\]
there is a canonical derived labelled execution diagram
\[
\mathbb A
\longleftarrow
\mathsf{Exec}^h_{\chi,\rho}
\longrightarrow
\mathbb Y.
\]
If the product
\[
\mathbb A\times\mathbb Y
\]
is formed in \(\mathcal C\), this diagram is equivalently a single map
\[
\mathsf{Exec}^h_{\chi,\rho}
\to
\mathbb A\times\mathbb Y.
\]
\end{proposition}

\begin{proof}
The two maps are the input projection
\[
I^h_{\chi,\rho}
\]
and downstream comparison
\[
\Lambda^h_{\chi,\rho}.
\]
The final statement is the product universal property in \(\mathcal C\).
\end{proof}

\section{Terminal trace and stateless cases}
\label{sec:derived-terminal-stateless-cases}

The terminal downstream behaviour is the identity map
\[
1_{\mathbb B}:\mathbb B\to\mathbb B.
\]

\begin{proposition}[Terminal trace case]
\label{prop:derived-terminal-trace-case}
For every derived Mealy morphism
\[
\chi=(\pi,\Omega):\mathbb X\to\mathbb A\times\mathbb B,
\]
there is a canonical equivalence
\[
\mathsf{Exec}^h_{\chi,1_{\mathbb B}}
=
\mathbb X\times_{\mathbb B}\mathbb B
\simeq
\mathbb X.
\]
Under this equivalence, the input projection is \(\pi\), and the downstream
comparison is the output cocycle
\[
\Omega:\mathbb X\to\mathbb B.
\]
\end{proposition}

\begin{proof}
The pullback of
\[
\Omega:\mathbb X\to\mathbb B
\]
along the identity map
\[
1_{\mathbb B}:\mathbb B\to\mathbb B
\]
is canonically equivalent to \(\mathbb X\). The two structural maps identify
with \(\pi\) and \(\Omega\).
\end{proof}

\begin{remark}[Trace interpretation]
\label{rem:derived-terminal-trace-interpretation}
The terminal downstream behaviour imposes no additional downstream dynamics.
Thus the terminal trace object of a derived Mealy morphism is simply its
derived execution object \(\mathbb X\), equipped with the output cocycle
\[
\Omega:\mathbb X\to\mathbb B.
\]
This is the derived version of the strict isomorphism
\[
P\times_{B_0}B_0\cong P.
\]
\end{remark}

\begin{definition}[Stateless derived Mealy morphism]
\label{def:stateless-derived-mealy-relative}
A derived Mealy morphism
\[
\chi:\mathbb X\to\mathbb A\times\mathbb B
\]
is \textbf{stateless} if its input map
\[
\pi:\mathbb X\to\mathbb A
\]
is equivalent in
\[
\mathcal C/\mathbb A
\]
to the identity map
\[
1_{\mathbb A}:\mathbb A\to\mathbb A.
\]
\end{definition}

A strict representative of a stateless derived Mealy morphism is supplied by a
derived internal functor
\[
F:\mathbb A\to\mathbb B.
\]
It determines
\[
\chi_F:\mathbb A\to\mathbb A\times\mathbb B
\]
classified by
\[
1_{\mathbb A}:\mathbb A\to\mathbb A,
\qquad
F:\mathbb A\to\mathbb B.
\]

\begin{proposition}[Stateless contextual restriction]
\label{prop:derived-stateless-contextual-restriction}
Let
\[
F:\mathbb A\to\mathbb B
\]
be a derived internal functor, and let
\[
\chi_F:\mathbb A\to\mathbb A\times\mathbb B
\]
be the associated stateless derived Mealy morphism. For every downstream
behaviour
\[
\rho:\mathbb Y\to\mathbb B,
\]
there is a canonical equivalence
\[
\mathsf{Exec}^h_{\chi_F,\rho}
\simeq
\mathbb A\times_{\mathbb B}\mathbb Y,
\]
where the pullback is taken over
\[
F:\mathbb A\to\mathbb B
\qquad
\text{and}
\qquad
\rho:\mathbb Y\to\mathbb B.
\]
Thus stateless contextual restriction is ordinary derived pullback of
downstream behaviours along \(F\).
\end{proposition}

\begin{proof}
For \(\chi_F\), the execution object is \(\mathbb A\) and the output map is
\(F\). Hence the contextual execution object is
\[
\mathbb A\times_{\mathbb B}\mathbb Y.
\]
\end{proof}

\begin{corollary}[Terminal stateless case]
\label{cor:derived-terminal-stateless-case}
For a stateless derived Mealy morphism classified by
\[
F:\mathbb A\to\mathbb B,
\]
one has
\[
\mathsf{Exec}^h_{\chi_F,1_{\mathbb B}}
\simeq
\mathbb A.
\]
Under the forward-opposite convention, this corresponds to the strict
identification
\[
\mathsf{Prog}_{F,\mathbf 1_{\mathbb B}}
\cong
\mathbb A^{\mathrm{op}}.
\]
\end{corollary}

\section{Functoriality in behaviours and Mealy morphisms}
\label{sec:derived-relative-functoriality}

Let
\[
\chi=(\pi,\Omega):\mathbb X\to\mathbb A\times\mathbb B
\]
be fixed.

\begin{proposition}[Functoriality in downstream behaviour]
\label{prop:derived-functoriality-downstream}
Let
\[
f:
(\rho:\mathbb Y\to\mathbb B)
\longrightarrow
(\rho':\mathbb Y'\to\mathbb B)
\]
be a morphism in
\[
\mathcal C/\mathbb B.
\]
Then there is an induced map
\[
\mathsf{Exec}^h_{\chi,\rho}
\longrightarrow
\mathsf{Exec}^h_{\chi,\rho'}
\]
given by
\[
1_{\mathbb X}\times_{\mathbb B}f.
\]
These maps are functorial in \(f\), and the resulting functor
\[
\mathcal C/\mathbb B\to\mathcal C/\mathbb A
\]
is precisely the contextual restriction functor
\[
\chi^\ast.
\]
\end{proposition}

\begin{proof}
The map
\[
f:\mathbb Y\to\mathbb Y'
\]
over \(\mathbb B\) induces a map of pullbacks
\[
\mathbb X\times_{\mathbb B}\mathbb Y
\to
\mathbb X\times_{\mathbb B}\mathbb Y'.
\]
Functoriality follows from functoriality of pullback in the slice. The map to
\(\mathbb A\) is obtained by composing with
\[
\pi:\mathbb X\to\mathbb A,
\]
which is exactly the definition of \(\chi^\ast\).
\end{proof}

Now let
\[
\chi:\mathbb X\to\mathbb A\times\mathbb B
\qquad
\text{and}
\qquad
\psi:\mathbb X'\to\mathbb A\times\mathbb B
\]
be derived Mealy morphisms.

\begin{definition}[Derived Mealy transformation]
\label{def:derived-mealy-transformation-relative}
A \textbf{derived Mealy transformation}
\[
\theta:\chi\to\psi
\]
is a morphism in the slice \(\infty\)-category
\[
\mathcal C/(\mathbb A\times\mathbb B).
\]
Thus it is represented by a map
\[
\theta:\mathbb X\to\mathbb X'
\]
compatible with the maps to \(\mathbb A\) and \(\mathbb B\).
\end{definition}

\begin{proposition}[Contextual functoriality of derived Mealy transformations]
\label{prop:derived-mealy-transformations-contextual}
Let
\[
\theta:\chi\to\psi
\]
be a derived Mealy transformation. For every downstream behaviour
\[
\rho:\mathbb Y\to\mathbb B,
\]
there is an induced map
\[
\theta_\rho:
\mathsf{Exec}^h_{\chi,\rho}
\to
\mathsf{Exec}^h_{\psi,\rho}
\]
given by
\[
\theta\times_{\mathbb B}1_{\mathbb Y}.
\]
It commutes with the input maps to \(\mathbb A\) and with the downstream
comparison maps to \(\mathbb Y\). The assignment
\[
\rho\longmapsto\theta_\rho
\]
is natural in \(\rho\).
\end{proposition}

\begin{proof}
Since \(\theta\) is a morphism over \(\mathbb B\), it induces a map of
pullbacks
\[
\mathbb X\times_{\mathbb B}\mathbb Y
\to
\mathbb X'\times_{\mathbb B}\mathbb Y.
\]
Compatibility with the input maps follows from the fact that \(\theta\) is
also a morphism over \(\mathbb A\). Compatibility with the downstream
projections is immediate from the definition. Naturality in \(\rho\) follows
from functoriality of pullback.
\end{proof}

\begin{remark}
\label{rem:higher-transformations-contextual}
Because transformations of derived Mealy morphisms live in an
\(\infty\)-categorical slice, they have higher homotopies. A path between
derived Mealy transformations induces a homotopy between the corresponding
contextual execution maps, and higher paths induce higher coherences.
\end{remark}

\section{Recovery of strict relative program categories}
\label{sec:strict-relative-recovery-derived}

We now compare the intrinsic derived construction with the strict relative
program categories of the previous chapter.

Let
\[
\Phi=(P,\delta,\omega):\mathbb A\rightsquigarrow\mathbb B
\]
be a strict Mealy morphism. It determines the backward action category
\[
P\rtimes\mathbb A
\]
and a point-set representative
\[
\mathbb A
\xleftarrow{\pi_P}
P\rtimes\mathbb A
\xrightarrow{\Omega_P}
\mathbb B.
\]
Here
\[
(P\rtimes\mathbb A)_0=P,
\qquad
(P\rtimes\mathbb A)_1=A_1\times_{t_A,A_0,p}P,
\]
and an arrow
\[
(a,x)
\]
has source \(\delta(a,x)\) and target \(x\).

Let
\[
Y=(Y\xrightarrow{r}B_0,\sigma)
\]
be a strict right \(\mathbb B\)-action. It determines the backward action
category
\[
Y\rtimes\mathbb B
\]
and a split internal discrete fibration
\[
\pi_Y:Y\rtimes\mathbb B\to\mathbb B.
\]

\begin{theorem}[Strict relative program categories as representatives]
\label{thm:strict-relative-programs-as-derived-representatives}
There is a canonical isomorphism of internal categories
\[
(P\rtimes\mathbb A)
\times_{\mathbb B}
(Y\rtimes\mathbb B)
\cong
\left(\mathsf{Prog}_{\Phi,Y}\right)^{\mathrm{op}},
\]
where the pullback on the left is the ordinary point-set pullback in
\(\mathbf{Cat}(\mathcal E)\), and \(\mathsf{Prog}_{\Phi,Y}\) is the
forward-oriented strict relative program category.

If the point-set cospan
\[
P\rtimes\mathbb A
\xrightarrow{\Omega_P}
\mathbb B
\xleftarrow{\pi_Y}
Y\rtimes\mathbb B
\]
is homotopy-pullback-ready, then this ordinary pullback represents the derived
execution object
\[
\mathsf{Exec}^h_{\chi_\Phi,\pi_Y}
\]
in
\[
\mathcal C.
\]
\end{theorem}

\begin{proof}
On objects, the pullback category has object object
\[
P\times_{q,B_0,r}Y,
\]
the strict relative state object.

An arrow in the pullback category consists of an arrow
\[
(a,x)
\]
of
\[
P\rtimes\mathbb A
\]
and an arrow
\[
(\beta,y)
\]
of
\[
Y\rtimes\mathbb B
\]
with the same image in \(\mathbb B\). Equality in the pullback forces
\[
\beta=\omega(a,x).
\]
Thus an arrow is uniquely of the form
\[
\bigl((a,x),(\omega(a,x),y)\bigr).
\]
Its source is
\[
\left(
\delta(a,x),
\sigma(\omega(a,x),y)
\right),
\]
and its target is
\[
(x,y).
\]
This is exactly the reverse of the forward strict relative transition
\[
(x,y)
\longrightarrow
\left(
\delta(a,x),
\sigma(\omega(a,x),y)
\right)
\]
in
\[
\mathsf{Prog}_{\Phi,Y}.
\]
The identity and composition maps agree after taking opposites, using the
Mealy unit and multiplication laws together with the action laws for \(Y\).

The final statement follows from the point-set pullback convention: whenever
the displayed cospan has been arranged to compute the homotopy pullback, its
ordinary pullback presents the corresponding pullback in \(\mathcal C\).
\end{proof}

\begin{corollary}[Strict restriction as derived contextual restriction]
\label{cor:strict-restriction-derived-contextual}
Under the representative comparison of
Theorem~\ref{thm:strict-relative-programs-as-derived-representatives}, the
derived contextual restriction functor
\[
\chi_\Phi^\ast:\mathcal C/\mathbb B\to\mathcal C/\mathbb A
\]
recovers the strict restriction of actions
\[
\Phi^\ast:\mathbf{Act}(\mathbb B)\to\mathbf{Act}(\mathbb A)
\]
after passing through split discrete-fibration representatives and applying
the chosen opposite convention.
\end{corollary}

\begin{remark}
\label{rem:opposite-convention-strict-derived}
The opposite in
\[
(P\rtimes\mathbb A)\times_{\mathbb B}(Y\rtimes\mathbb B)
\cong
\left(\mathsf{Prog}_{\Phi,Y}\right)^{\mathrm{op}}
\]
is not a defect. It records the same contravariant input-processing
orientation already present in the Mealy structure cell
\[
\Phi:P\circ A\Rightarrow B\circ P.
\]
The intrinsic derived execution object keeps this contravariant orientation.
The forward program category is obtained by reversing it.
\end{remark}

\section{Derived proarrow programs over contextual executions}
\label{sec:derived-proarrow-programs-contextual}

Let
\[
\chi:\mathbb A\rightsquigarrow^h\mathbb B
\]
be a derived Mealy morphism, let
\[
\rho:\mathbb Y\to\mathbb B
\]
be a downstream behaviour, and write
\[
\mathbb Z
:=
\mathsf{Exec}^h_{\chi,\rho}
=
\mathbb X\times_{\mathbb B}\mathbb Y.
\]

\begin{definition}[Derived contextual program]
\label{def:derived-contextual-program}
A \textbf{derived contextual program} on
\[
\mathbb Z
\]
is an endocorrespondence
\[
R:\mathbb Z\nrightarrow\mathbb Z,
\]
that is, a span
\[
\mathbb Z
\xleftarrow{\ell}
\mathbb R
\xrightarrow{r}
\mathbb Z
\]
in
\[
\mathcal C.
\]
The \(\infty\)-category of derived contextual programs on \(\mathbb Z\) is
\[
\operatorname{Corr}(\mathcal C)(\mathbb Z,\mathbb Z).
\]
\end{definition}

\begin{remark}
\label{rem:no-automatic-primitive-span-derived}
In the strict relative chapter, the primitive execution span is automatically
the arrow span of the program category
\[
\mathsf{Prog}_{\Phi,Y}.
\]
In the intrinsic derived setting, the contextual execution object
\[
\mathbb Z
\]
does not by itself determine a primitive endoprogram. A primitive program must
be supplied either by a point-set representative, by an additional derived
endocorrespondence, or by a derived execution monoid.
\end{remark}

\begin{definition}[Derived primitive program and subprogram]
\label{def:derived-primitive-program-subprogram}
A \textbf{derived primitive program} on
\[
\mathbb Z
\]
is a chosen endocorrespondence
\[
\mathbf m:\mathbb Z\nrightarrow\mathbb Z.
\]
A \textbf{subprogram} of \(\mathbf m\) is a morphism of endocorrespondences
\[
U\to\mathbf m.
\]
When the relevant correspondence category is posetal over \(\mathbf m\), such
a subprogram may be regarded as a subobject of the apex of \(\mathbf m\).
\end{definition}

\section{Derived predicate semantics}
\label{sec:derived-predicate-semantics}

The strict dynamic proarrow logic of the previous chapter was expressed using
ordinary subobject lattices and regular images. In the derived setting this
must be separated into two levels: representative-level predicate semantics
and intrinsic derived predicate semantics.

\subsection{Representative-level predicate semantics}
\label{subsec:representative-level-predicate-semantics}

Given a point-set representative of a derived contextual execution object,
one may apply the strict predicate semantics to the corresponding point-set
state object and execution spans. This gives formulas involving
\[
\mathrm{Sub}(S),
\qquad
\exists_f,
\qquad
\forall_f.
\]
This semantics is useful for recovering the strict relative chapter, but it is
not by itself an invariant derived definition unless the chosen representative
and predicate semantics are known to respect weak equivalences.

\subsection{Intrinsic derived predicate semantics}
\label{subsec:intrinsic-derived-predicate-semantics}

For intrinsic derived proarrow logic we assume an abstract predicate theory on
\[
\mathcal C.
\]

\begin{definition}[Derived regular predicate theory]
\label{def:derived-regular-predicate-theory}
A \textbf{derived regular predicate theory} on \(\mathcal C\) consists of:
\begin{enumerate}
\item for every object \(\mathbb Z\in\mathcal C\), a poset
\[
\mathrm{Pred}(\mathbb Z)
\]
of predicates on \(\mathbb Z\);
\item for every morphism
\[
f:\mathbb X\to\mathbb Y
\]
a pullback operation
\[
f^\ast:\mathrm{Pred}(\mathbb Y)\to\mathrm{Pred}(\mathbb X);
\]
\item for every such \(f\), a left adjoint
\[
\Sigma_f:\mathrm{Pred}(\mathbb X)\to\mathrm{Pred}(\mathbb Y)
\]
to \(f^\ast\);
\item Beck--Chevalley equivalences for \(\Sigma\) along pullback squares in
\(\mathcal C\).
\end{enumerate}
Thus
\[
\Sigma_f\dashv f^\ast.
\]
\end{definition}

\begin{definition}[Derived first-order Heyting predicate theory]
\label{def:derived-heyting-predicate-theory}
A \textbf{derived first-order Heyting predicate theory} on \(\mathcal C\) is a
derived regular predicate theory such that:
\begin{enumerate}
\item each fibre \(\mathrm{Pred}(\mathbb Z)\) has finite meets, finite joins,
and Heyting implication;
\item every pullback functor
\[
f^\ast:\mathrm{Pred}(\mathbb Y)\to\mathrm{Pred}(\mathbb X)
\]
has a right adjoint
\[
\Pi_f:\mathrm{Pred}(\mathbb X)\to\mathrm{Pred}(\mathbb Y);
\]
\item the right adjoints \(\Pi_f\) satisfy Beck--Chevalley for pullback
squares.
\end{enumerate}
Thus, for every \(f\), there is an adjoint triple
\[
\Sigma_f\dashv f^\ast\dashv \Pi_f.
\]
\end{definition}

\begin{remark}
\label{rem:predicate-theory-not-automatic}
We do not assert that
\[
\mathcal C=\mathbf{Cat}(\mathcal E)_{JT,\infty}
\]
automatically carries the full predicate structure above. The logical layer of
this chapter is conditional on the chosen derived predicate theory. In
applications, such a predicate theory may come from an ambient
\(\infty\)-topos, a stack semantics, or a suitably homotopy-invariant
point-set predicate doctrine.
\end{remark}

\section{Diamonds, boxes, and tests for derived proarrows}
\label{sec:derived-diamonds-boxes-tests}

Assume first that \(\mathcal C\) is equipped with a derived regular predicate
theory.

\begin{definition}[Derived diamond of a correspondence]
\label{def:derived-diamond-correspondence}
Let
\[
R=
\left(
\mathbb X
\xleftarrow{\ell}
\mathbb R
\xrightarrow{r}
\mathbb Y
\right)
\]
be a correspondence in \(\mathcal C\). Its \textbf{diamond modality} is
\[
\langle R\rangle:
\mathrm{Pred}(\mathbb Y)\to\mathrm{Pred}(\mathbb X)
\]
defined by
\[
\langle R\rangle\varphi
:=
\Sigma_\ell(r^\ast\varphi).
\]
\end{definition}

\begin{proposition}[Functoriality of derived diamonds]
\label{prop:derived-diamond-functoriality}
For correspondences
\[
R:\mathbb X\nrightarrow\mathbb Y
\qquad
\text{and}
\qquad
S:\mathbb Y\nrightarrow\mathbb Z,
\]
one has
\[
\langle 1_{\mathbb X}\rangle=\mathrm{id}_{\mathrm{Pred}(\mathbb X)}
\]
and
\[
\langle R;S\rangle
=
\langle R\rangle\circ\langle S\rangle.
\]
\end{proposition}

\begin{proof}
The identity law follows from the unit correspondence. For composition, write
\[
R=(\mathbb X\xleftarrow{\ell}\mathbb R\xrightarrow{r}\mathbb Y),
\qquad
S=(\mathbb Y\xleftarrow{m}\mathbb S\xrightarrow{n}\mathbb Z).
\]
Their composite has apex
\[
\mathbb R\times_{\mathbb Y}\mathbb S.
\]
Thus
\[
\langle R;S\rangle\varphi
=
\Sigma_{\ell\pi_{\mathbb R}}
\bigl((n\pi_{\mathbb S})^\ast\varphi\bigr).
\]
Using functoriality of \(\Sigma\) and Beck--Chevalley for the pullback square
defining \(\mathbb R\times_{\mathbb Y}\mathbb S\), this becomes
\[
\Sigma_\ell r^\ast\Sigma_m n^\ast\varphi
=
\langle R\rangle(\langle S\rangle\varphi).
\]
\end{proof}

Now assume that the predicate theory is derived first-order Heyting.

\begin{definition}[Derived box of a correspondence]
\label{def:derived-box-correspondence}
Let
\[
R=
\left(
\mathbb X
\xleftarrow{\ell}
\mathbb R
\xrightarrow{r}
\mathbb Y
\right)
\]
be a correspondence. Its \textbf{box modality} is
\[
[R]:
\mathrm{Pred}(\mathbb Y)\to\mathrm{Pred}(\mathbb X)
\]
defined by
\[
[R]\varphi
:=
\Pi_\ell(r^\ast\varphi).
\]
\end{definition}

\begin{proposition}[Functoriality of derived boxes]
\label{prop:derived-box-functoriality}
For correspondences
\[
R:\mathbb X\nrightarrow\mathbb Y
\qquad
\text{and}
\qquad
S:\mathbb Y\nrightarrow\mathbb Z,
\]
one has
\[
[1_{\mathbb X}]=\mathrm{id}_{\mathrm{Pred}(\mathbb X)}
\]
and
\[
[R;S]=[R]\circ[S].
\]
\end{proposition}

\begin{proof}
The proof is the same as the proof of
Proposition~\ref{prop:derived-diamond-functoriality}, using the right
adjoints
\[
\Pi_f
\]
and their Beck--Chevalley equivalences instead of the left adjoints
\[
\Sigma_f.
\]
\end{proof}

\begin{definition}[Derived test]
\label{def:derived-test}
Let
\[
\theta\in\mathrm{Pred}(\mathbb Z)
\]
be represented by a predicate inclusion
\[
\Theta\hookrightarrow\mathbb Z.
\]
The corresponding \textbf{test correspondence} is
\[
\theta?
=
\left(
\mathbb Z
\xleftarrow{\theta}
\Theta
\xrightarrow{\theta}
\mathbb Z
\right).
\]
\end{definition}

\begin{proposition}[Diamond law for derived tests]
\label{prop:derived-test-diamond-law}
For every predicate
\[
\varphi\in\mathrm{Pred}(\mathbb Z),
\]
one has
\[
\langle\theta?\rangle\varphi
=
\theta\wedge\varphi.
\]
\end{proposition}

\begin{proof}
By definition,
\[
\langle\theta?\rangle\varphi
=
\Sigma_\theta(\theta^\ast\varphi).
\]
For a predicate inclusion, the dependent sum along \(\theta\) identifies
predicates on \(\Theta\) with predicates on \(\mathbb Z\) lying below
\(\theta\). Hence the image of \(\theta^\ast\varphi\) is
\[
\theta\wedge\varphi.
\]
\end{proof}

\begin{proposition}[Box law for derived tests]
\label{prop:derived-test-box-law}
In a derived first-order Heyting predicate theory,
\[
[\theta?]\varphi
=
\theta\Rightarrow\varphi.
\]
\end{proposition}

\begin{proof}
By definition,
\[
[\theta?]\varphi
=
\Pi_\theta(\theta^\ast\varphi).
\]
For any predicate \(\psi\in\mathrm{Pred}(\mathbb Z)\),
\[
\psi\le \Pi_\theta(\theta^\ast\varphi)
\]
if and only if
\[
\theta^\ast\psi\le \theta^\ast\varphi.
\]
Since \(\theta\) is a predicate inclusion, this is equivalent to
\[
\psi\wedge\theta\le \varphi.
\]
By the defining property of Heyting implication, this is equivalent to
\[
\psi\le \theta\Rightarrow\varphi.
\]
\end{proof}

\section{Contextual relative modalities}
\label{sec:contextual-relative-modalities-derived}

Let
\[
\chi:\mathbb A\rightsquigarrow^h\mathbb B
\]
be a derived Mealy morphism, let
\[
\rho:\mathbb Y\to\mathbb B
\]
be a downstream behaviour, and write
\[
\mathbb Z
:=
\mathsf{Exec}^h_{\chi,\rho}.
\]

\begin{definition}[Derived contextual modality]
\label{def:derived-contextual-modality}
Let
\[
U:\mathbb Z\nrightarrow\mathbb Z
\]
be a derived contextual program. Its diamond is
\[
\langle U\rangle:
\mathrm{Pred}(\mathbb Z)\to\mathrm{Pred}(\mathbb Z),
\]
given by
\[
\langle U\rangle\varphi
=
\Sigma_\ell(r^\ast\varphi)
\]
for
\[
U=(\mathbb Z\xleftarrow{\ell}\mathbb U\xrightarrow{r}\mathbb Z).
\]
If the predicate theory is first-order Heyting, its box is
\[
[U]\varphi
=
\Pi_\ell(r^\ast\varphi).
\]
\end{definition}

\begin{remark}[Interpretation]
\label{rem:contextual-modality-interpretation}
The predicate
\[
\langle U\rangle\varphi
\]
means that there exists a \(U\)-execution from the current derived contextual
state to a derived contextual state satisfying \(\varphi\). The predicate
\[
[U]\varphi
\]
means that every \(U\)-execution from the current derived contextual state
lands in a state satisfying \(\varphi\).
\end{remark}

\begin{definition}[Label-restricted derived subprogram]
\label{def:label-restricted-derived-subprogram}
Suppose a derived contextual program
\[
U=(\mathbb Z\xleftarrow{\ell}\mathbb U\xrightarrow{r}\mathbb Z)
\]
is equipped with a label map
\[
\lambda:\mathbb U\to\mathbb L.
\]
For a predicate
\[
R\in\mathrm{Pred}(\mathbb L),
\]
the \textbf{restriction of \(U\) by \(R\)} is the subprogram obtained by
pulling \(R\) back along \(\lambda\):
\[
U|_R
\hookrightarrow
U.
\]
Its modality is denoted
\[
\langle U|_R\rangle.
\]
\end{definition}

\begin{remark}
\label{rem:label-maps-not-automatic-derived}
In the intrinsic derived setting, label maps on the apex of a program
correspondence are additional structure. They may arise from a point-set
representative, from a chosen primitive correspondence, or from a richer
universal construction. This is the derived replacement of the strict
restrictions by input, raw-output, and downstream-transition predicates.
\end{remark}

\section{Finite choice of derived subprograms}
\label{sec:finite-choice-derived-subprograms}

Assume that the relevant predicate fibres and subprogram fibres have finite
joins stable under pullback.

Let
\[
\mathbf m:\mathbb Z\nrightarrow\mathbb Z
\]
be a chosen derived primitive program, and let
\[
U,V\to\mathbf m
\]
be subprograms. Their finite choice is written
\[
U\vee V\to\mathbf m.
\]

\begin{proposition}[Finite choice law]
\label{prop:derived-finite-choice-law}
For every predicate
\[
\varphi\in\mathrm{Pred}(\mathbb Z),
\]
one has
\[
\langle U\vee V\rangle\varphi
=
\langle U\rangle\varphi
\vee
\langle V\rangle\varphi.
\]
\end{proposition}

\begin{proof}
Regard \(U\) and \(V\) as predicates on the apex of \(\mathbf m\). If
\[
\mathbf m=(\mathbb Z\xleftarrow{s}\mathbb M\xrightarrow{t}\mathbb Z),
\]
then
\[
\langle U\rangle\varphi
=
\Sigma_s(U\wedge t^\ast\varphi).
\]
Hence
\[
\langle U\vee V\rangle\varphi
=
\Sigma_s((U\vee V)\wedge t^\ast\varphi).
\]
By distributivity of finite joins over finite meets,
\[
(U\vee V)\wedge t^\ast\varphi
=
(U\wedge t^\ast\varphi)
\vee
(V\wedge t^\ast\varphi).
\]
Since \(\Sigma_s\) is a left adjoint, it preserves finite joins. Therefore
\[
\langle U\vee V\rangle\varphi
=
\langle U\rangle\varphi
\vee
\langle V\rangle\varphi.
\]
\end{proof}

\begin{remark}
\label{rem:derived-choice-fixed-primitive}
This is finite choice of subprograms inside a fixed primitive program
\[
\mathbf m.
\]
It should not be confused with arbitrary coproducts of unrelated
correspondences unless the relevant coproduct hypotheses have been imposed.
\end{remark}

\section{Path star in the derived correspondence category}
\label{sec:path-star-derived-corr}

Let
\[
\mathbb Z\in\mathcal C
\]
and let
\[
U:\mathbb Z\nrightarrow\mathbb Z
\]
be an endocorrespondence.

Assume that the hom-\(\infty\)-category
\[
\operatorname{Corr}(\mathcal C)(\mathbb Z,\mathbb Z)
\]
has countable coproducts and that composition preserves these coproducts
separately in each variable.

Define powers by
\[
U^0:=1_{\mathbb Z},
\qquad
U^{n+1}:=U;U^n.
\]

\begin{definition}[Derived finite-path star]
\label{def:derived-path-star}
The \textbf{finite-path star} of \(U\) is
\[
U^\star_{\mathrm{path}}
:=
\coprod_{n\ge0}U^n.
\]
\end{definition}

\begin{proposition}[Derived path-star equation]
\label{prop:derived-path-star-equation}
There is a canonical equivalence of endocorrespondences
\[
U^\star_{\mathrm{path}}
\simeq
1_{\mathbb Z}\oplus U;U^\star_{\mathrm{path}}.
\]
\end{proposition}

\begin{proof}
By definition,
\[
U^\star_{\mathrm{path}}
=
\coprod_{n\ge0}U^n.
\]
Separating the \(0\)-summand gives
\[
\coprod_{n\ge0}U^n
\simeq
U^0\oplus\coprod_{n\ge0}U^{n+1}.
\]
Since
\[
U^0=1_{\mathbb Z}
\]
and
\[
U^{n+1}=U;U^n,
\]
and since composition with \(U\) preserves the relevant countable coproducts,
\[
\coprod_{n\ge0}U^{n+1}
\simeq
U;\left(\coprod_{n\ge0}U^n\right)
=
U;U^\star_{\mathrm{path}}.
\]
\end{proof}

\begin{theorem}[Derived path star as free monoid]
\label{thm:derived-path-star-free-monoid}
Under the hypotheses above,
\[
U^\star_{\mathrm{path}}
=
\coprod_{n\ge0}U^n
\]
is the free monoid on \(U\) in the monoidal \(\infty\)-category
\[
\operatorname{Corr}(\mathcal C)(\mathbb Z,\mathbb Z).
\]
\end{theorem}

\begin{proof}
The monoidal product is correspondence composition. Because this product
preserves countable coproducts separately in each variable, the usual tensor
algebra construction applies.

The unit is the inclusion of the \(0\)-summand
\[
1_{\mathbb Z}=U^0\to U^\star_{\mathrm{path}}.
\]
Multiplication is induced summandwise by the canonical equivalences
\[
U^m;U^n\simeq U^{m+n}
\]
followed by the coproduct inclusion
\[
U^{m+n}\to \coprod_{k\ge0}U^k.
\]
Associativity and the unit laws follow from associativity and unitality of
correspondence composition together with associativity and unitality of
addition of natural numbers.

Given any monoid object \(M\) and a map \(U\to M\), the induced maps
\[
U^n\to M
\]
are defined by iterated multiplication in \(M\). The coproduct universal
property gives a unique map
\[
U^\star_{\mathrm{path}}\to M
\]
compatible with these maps, and the compatibility is exactly the monoid
morphism condition.
\end{proof}

\section{Fixed-point semantics for derived path star}
\label{sec:derived-fixed-point-semantics}

Assume now that the derived predicate theory has countable joins and countable
meets in predicate fibres, and that the modalities considered below preserve
the relevant countable joins or meets.

Let
\[
U:\mathbb Z\nrightarrow\mathbb Z
\]
be an endocorrespondence and let
\[
\varphi\in\mathrm{Pred}(\mathbb Z).
\]

\subsection{Diamond fixed points}
\label{subsec:derived-diamond-fixed-points}

Define
\[
F^\Diamond_{U,\varphi}:\mathrm{Pred}(\mathbb Z)\to
\mathrm{Pred}(\mathbb Z)
\]
by
\[
F^\Diamond_{U,\varphi}(X)
:=
\varphi\vee\langle U\rangle X.
\]
Assume that
\[
\langle U\rangle
\]
preserves countable joins. Define
\[
\mu F^\Diamond_{U,\varphi}
:=
\bigvee_{n\ge0}(F^\Diamond_{U,\varphi})^n(\bot).
\]

\begin{lemma}[Derived diamond approximants]
\label{lem:derived-diamond-approximants}
For every \(n\ge0\),
\[
(F^\Diamond_{U,\varphi})^n(\bot)
=
\bigvee_{0\le k<n}\langle U^k\rangle\varphi.
\]
\end{lemma}

\begin{proof}
The case \(n=0\) is the empty join. The induction step uses
\[
F^\Diamond_{U,\varphi}(X)=\varphi\vee\langle U\rangle X,
\]
preservation of joins by \(\langle U\rangle\), and functoriality of diamonds:
\[
\langle U\rangle\langle U^k\rangle
=
\langle U;U^k\rangle
=
\langle U^{k+1}\rangle.
\]
\end{proof}

\begin{theorem}[Derived diamond least fixed point]
\label{thm:derived-diamond-lfp}
One has
\[
\mu X.\bigl(\varphi\vee\langle U\rangle X\bigr)
=
\bigvee_{n\ge0}\langle U^n\rangle\varphi.
\]
This predicate is the least pre-fixed point of
\[
X\longmapsto \varphi\vee\langle U\rangle X.
\]
\end{theorem}

\begin{proof}
The displayed formula follows from
Lemma~\ref{lem:derived-diamond-approximants}. Let
\[
\chi:=
\bigvee_{n\ge0}\langle U^n\rangle\varphi.
\]
Then
\[
\begin{aligned}
\varphi\vee\langle U\rangle\chi
&=
\langle U^0\rangle\varphi
\vee
\bigvee_{n\ge0}\langle U\rangle\langle U^n\rangle\varphi
\\
&=
\langle U^0\rangle\varphi
\vee
\bigvee_{n\ge0}\langle U^{n+1}\rangle\varphi
\\
&=
\chi.
\end{aligned}
\]
If
\[
\varphi\vee\langle U\rangle\psi\le\psi,
\]
then \(\varphi\le\psi\) and \(\langle U\rangle\psi\le\psi\). By induction,
\[
\langle U^n\rangle\varphi\le\psi
\]
for all \(n\), hence
\[
\chi\le\psi.
\]
Thus \(\chi\) is the least pre-fixed point.
\end{proof}

\subsection{Box fixed points}
\label{subsec:derived-box-fixed-points}

Assume the predicate theory is first-order Heyting and that
\[
[U]
\]
preserves countable meets. Define
\[
F^\Box_{U,\varphi}(X)
:=
\varphi\wedge[U]X.
\]
Set
\[
\nu F^\Box_{U,\varphi}
:=
\bigwedge_{n\ge0}(F^\Box_{U,\varphi})^n(\top).
\]

\begin{lemma}[Derived box approximants]
\label{lem:derived-box-approximants}
For every \(n\ge0\),
\[
(F^\Box_{U,\varphi})^n(\top)
=
\bigwedge_{0\le k<n}[U^k]\varphi.
\]
\end{lemma}

\begin{proof}
The proof is dual to Lemma~\ref{lem:derived-diamond-approximants}, using
preservation of meets by \([U]\) and functoriality of boxes:
\[
[U][U^k]=[U;U^k]=[U^{k+1}].
\]
\end{proof}

\begin{theorem}[Derived box greatest fixed point]
\label{thm:derived-box-gfp}
One has
\[
\nu X.\bigl(\varphi\wedge[U]X\bigr)
=
\bigwedge_{n\ge0}[U^n]\varphi.
\]
This predicate is the greatest post-fixed point of
\[
X\longmapsto \varphi\wedge[U]X.
\]
\end{theorem}

\begin{proof}
Let
\[
\chi:=
\bigwedge_{n\ge0}[U^n]\varphi.
\]
By Lemma~\ref{lem:derived-box-approximants}, this is the displayed
\(\nu\)-approximant. Moreover,
\[
\begin{aligned}
\varphi\wedge[U]\chi
&=
[U^0]\varphi
\wedge
\bigwedge_{n\ge0}[U][U^n]\varphi
\\
&=
[U^0]\varphi
\wedge
\bigwedge_{n\ge0}[U^{n+1}]\varphi
\\
&=
\chi.
\end{aligned}
\]
If
\[
\psi\le \varphi\wedge[U]\psi,
\]
then \(\psi\le\varphi\) and \(\psi\le[U]\psi\). By induction,
\[
\psi\le[U^n]\varphi
\]
for all \(n\), hence
\[
\psi\le\chi.
\]
Thus \(\chi\) is the greatest post-fixed point.
\end{proof}

\begin{theorem}[Derived path/fixed-point comparison]
\label{thm:derived-path-fixed-point-comparison}
Assume the modalities of countable coproducts satisfy
\[
\left\langle\coprod_{n\ge0}R_n\right\rangle\varphi
=
\bigvee_{n\ge0}\langle R_n\rangle\varphi
\]
and, in the box case,
\[
\left[\coprod_{n\ge0}R_n\right]\varphi
=
\bigwedge_{n\ge0}[R_n]\varphi.
\]
Then
\[
\langle U^\star_{\mathrm{path}}\rangle\varphi
=
\mu X.\bigl(\varphi\vee\langle U\rangle X\bigr),
\]
and, when boxes exist,
\[
[U^\star_{\mathrm{path}}]\varphi
=
\nu X.\bigl(\varphi\wedge[U]X\bigr).
\]
\end{theorem}

\begin{proof}
By definition,
\[
U^\star_{\mathrm{path}}
=
\coprod_{n\ge0}U^n.
\]
Using the coproduct modality formula,
\[
\langle U^\star_{\mathrm{path}}\rangle\varphi
=
\bigvee_{n\ge0}\langle U^n\rangle\varphi.
\]
By Theorem~\ref{thm:derived-diamond-lfp}, this is
\[
\mu X.\bigl(\varphi\vee\langle U\rangle X\bigr).
\]
The box statement is the same, using the coproduct formula for boxes and
Theorem~\ref{thm:derived-box-gfp}.
\end{proof}

\begin{corollary}[Derived induction and coinduction]
\label{cor:derived-induction-coinduction}
If
\[
\varphi\vee\langle U\rangle\psi\le\psi,
\]
then
\[
\mu X.\bigl(\varphi\vee\langle U\rangle X\bigr)\le\psi.
\]
If boxes exist and
\[
\psi\le\varphi\wedge[U]\psi,
\]
then
\[
\psi\le\nu X.\bigl(\varphi\wedge[U]X\bigr).
\]
\end{corollary}

\section{Derived folded trace semantics}
\label{sec:derived-folded-trace-semantics}

Let
\[
\mathbb Z=\mathsf{Exec}^h_{\chi,\rho}.
\]
A folded trace construction requires more than the contextual execution object:
it requires a chosen execution monoid.

\begin{definition}[Derived execution monoid]
\label{def:derived-execution-monoid}
A \textbf{derived execution monoid} on
\[
\mathbb Z
\]
is a monoid object
\[
\mathbf m:\mathbb Z\nrightarrow\mathbb Z
\]
in the monoidal \(\infty\)-category
\[
\operatorname{Corr}(\mathcal C)(\mathbb Z,\mathbb Z).
\]
\end{definition}

Let
\[
i:U\to\mathbf m
\]
be a subprogram of a derived execution monoid.

\begin{definition}[Derived folded trace map]
\label{def:derived-folded-trace-map}
The \textbf{derived folded trace map} of
\[
i:U\to\mathbf m
\]
is the unique monoid morphism
\[
\operatorname{fold}_U:
U^\star_{\mathrm{path}}\to\mathbf m
\]
whose restriction to the \(1\)-step summand \(U\) is \(i\).
\end{definition}

\begin{proposition}[Meaning of the folded trace map]
\label{prop:meaning-derived-folded-trace}
The map
\[
\operatorname{fold}_U:
U^\star_{\mathrm{path}}\to\mathbf m
\]
sends a finite path of \(U\)-steps to its composite execution in the derived
execution monoid
\[
\mathbf m.
\]
\end{proposition}

\begin{proof}
The path star
\[
U^\star_{\mathrm{path}}
\]
is the free monoid on \(U\). A map
\[
i:U\to\mathbf m
\]
into a monoid object \(\mathbf m\) therefore extends uniquely to a monoid
morphism
\[
U^\star_{\mathrm{path}}\to\mathbf m.
\]
On the \(n\)-step summand \(U^n\), this map is given by iterated
multiplication in \(\mathbf m\). Hence it folds a finite path into its
composite execution.
\end{proof}

\begin{remark}[Strict comparison]
\label{rem:strict-comparison-folded-derived}
For a strict Mealy morphism and a strict downstream action, the primitive
execution span of the forward program category
\[
\mathsf{Prog}_{\Phi,Y}
\]
is a monoid object in the appropriate endospan category. Its derived
representative supplies an execution monoid. The derived folded trace map then
recovers the strict folded trace map after applying the point-set comparison
and the chosen opposite convention.
\end{remark}

\section{Representative-level primitive semantics}
\label{sec:representative-level-primitive-semantics}

We spell out how the familiar strict formulas are recovered from
representatives.

Let
\[
\Phi=(P,\delta,\omega):\mathbb A\rightsquigarrow\mathbb B
\]
be a strict Mealy morphism and let
\[
Y=(Y\xrightarrow{r}B_0,\sigma)
\]
be a strict right \(\mathbb B\)-action. The forward relative state object is
\[
S_{\Phi,Y}:=P\times_{q,B_0,r}Y.
\]
The forward primitive execution object is
\[
M_{\Phi,Y}
:=
A_1\times_{t_A,A_0,p\pi_P}S_{\Phi,Y}.
\]
A generalized element is a triple
\[
(a,x,y)
\]
with
\[
t_A(a)=p(x),
\qquad
q(x)=r(y).
\]
The source and target maps are
\[
s_{\Phi,Y}(a,x,y)=(x,y),
\]
and
\[
t_{\Phi,Y}(a,x,y)
=
\left(
\delta(a,x),
\sigma(\omega(a,x),y)
\right).
\]

\begin{definition}[Representative primitive diamond]
\label{def:representative-primitive-diamond}
In a regular point-set predicate semantics, the representative primitive
diamond is
\[
\langle\mathbf m_{\Phi,Y}\rangle\varphi
:=
\exists_{s_{\Phi,Y}}\bigl(t_{\Phi,Y}^\ast\varphi\bigr).
\]
\end{definition}

\begin{proposition}[External reading of representative primitive diamonds]
\label{prop:representative-primitive-diamond-reading}
In \(\mathbf{Set}\),
\[
(x,y)\models\langle\mathbf m_{\Phi,Y}\rangle\varphi
\]
if and only if there exists an admissible input arrow
\[
a:u\to p(x)
\]
such that
\[
\left(
\delta(a,x),
\sigma(\omega(a,x),y)
\right)
\models\varphi.
\]
\end{proposition}

\begin{proof}
This is the elementwise interpretation of
\[
\exists_{s_{\Phi,Y}}\bigl(t_{\Phi,Y}^\ast\varphi\bigr).
\]
A witness is a primitive execution
\[
(a,x,y)\in M_{\Phi,Y}
\]
with source \((x,y)\) and target
\[
\left(
\delta(a,x),
\sigma(\omega(a,x),y)
\right).
\]
\end{proof}

\begin{remark}
\label{rem:representative-semantics-not-intrinsic}
The modality in
Definition~\ref{def:representative-primitive-diamond} is a point-set
representative semantics. It is the correct way to recover the strict relative
program logic, but it is not the intrinsic derived definition unless the
chosen representative predicate semantics is known to be homotopy invariant.
\end{remark}

\section{Classical and relational special cases}
\label{sec:derived-relative-classical-special-cases}

\begin{example}[Classical Mealy cascades]
\label{ex:derived-classical-mealy-cascades}
Let
\[
\mathcal E=\mathbf{Set}
\]
and suppose that \(\mathbb A\) and \(\mathbb B\) each have one object. Then
they are monoids, say
\[
M
\qquad
\text{and}
\qquad
N.
\]
A strict Mealy morphism consists of a set \(P\) and functions
\[
\delta:M\times P\to P,
\qquad
\omega:M\times P\to N
\]
satisfying the unit and cocycle laws.

A right \(N\)-action is a set \(Y\) with an action
\[
N\times Y\to Y.
\]
The contextual state object is
\[
P\times Y,
\]
and the induced \(M\)-action is
\[
m\cdot(x,y)
=
\left(
\delta(m,x),
\omega(m,x)\cdot y
\right).
\]
Thus the strict representative of the derived contextual execution object is
the ordinary automata-theoretic cascade of a Mealy transducer followed by an
\(N\)-automaton.
\end{example}

\begin{example}[Relational dynamic logic]
\label{ex:derived-relative-relational-pdl}
In \(\mathbf{Set}\), after quotienting spans by their extensional binary
relations, a span
\[
S\xleftarrow{\ell}R\xrightarrow{r}S
\]
determines a relation
\[
s\,\overline R\,s'
\]
if and only if there exists \(\rho\in R\) such that
\[
\ell(\rho)=s,
\qquad
r(\rho)=s'.
\]
The diamond becomes the usual relational dynamic-logic modality:
\[
s\models\langle R\rangle\varphi
\]
if and only if there exists \(s'\) such that
\[
s\,\overline R\,s'
\]
and
\[
s'\models\varphi.
\]
For a relative Mealy system, this relation lives on the contextual state space
\[
P\times_{B_0}Y.
\]
\end{example}

\begin{example}[Relational Kleene algebra with tests]
\label{ex:derived-relative-relational-kat}
In \(\mathbf{Set}\), after passing from spans to extensional relations,
endoprograms on a state set \(S\) form the relation algebra
\[
\mathcal P(S\times S).
\]
Sequential composition is relational composition, finite choice is union,
tests are subidentity relations, and path star is reflexive-transitive
closure. For a relative Mealy system, this gives a Kleene algebra with tests
on the coupled state object
\[
P\times_{B_0}Y.
\]
\end{example}

\section{Main consolidation}
\label{sec:derived-relative-main-consolidation}

\begin{theorem}[Derived contextual Mealy proarrow logic]
\label{thm:derived-contextual-mealy-proarrow-logic}
Let
\[
\mathcal E
\]
be a Grothendieck topos, let
\[
\mathcal M=\mathbf{Cat}(\mathcal E)_{JT},
\]
and let
\[
\mathcal C=\mathcal M_\infty
\]
be the Joyal--Tierney localized \(\infty\)-category of internal categories.

\begin{enumerate}
\item A derived Mealy morphism from \(\mathbb A\) to \(\mathbb B\) is a
derived correspondence
\[
\chi:\mathbb X\to\mathbb A\times\mathbb B.
\]

\item A derived downstream \(\mathbb B\)-behaviour is an object
\[
\rho:\mathbb Y\to\mathbb B
\]
of the slice
\[
\mathcal C/\mathbb B.
\]

\item Every derived Mealy morphism
\[
\chi=(\pi,\Omega):\mathbb X\to\mathbb A\times\mathbb B
\]
induces a contextual restriction functor
\[
\chi^\ast:\mathcal C/\mathbb B\to\mathcal C/\mathbb A
\]
given by
\[
\rho:\mathbb Y\to\mathbb B
\longmapsto
\left(
\mathbb X\times_{\mathbb B}\mathbb Y\to\mathbb A
\right).
\]

\item The derived relative execution object in context \(\rho\) is
\[
\mathsf{Exec}^h_{\chi,\rho}
=
\mathbb X\times_{\mathbb B}\mathbb Y.
\]

\item It carries canonical structural maps
\[
\mathsf{Exec}^h_{\chi,\rho}\to\mathbb A
\qquad
\text{and}
\qquad
\mathsf{Exec}^h_{\chi,\rho}\to\mathbb Y.
\]

\item The terminal downstream behaviour gives
\[
\mathsf{Exec}^h_{\chi,1_{\mathbb B}}
\simeq
\mathbb X.
\]

\item If \(\chi\) is stateless and classified by a derived functor
\[
F:\mathbb A\to\mathbb B,
\]
then
\[
\mathsf{Exec}^h_{\chi,\rho}
\simeq
\mathbb A\times_{\mathbb B}\mathbb Y.
\]

\item Strict Mealy morphisms and strict downstream actions give point-set
representatives
\[
(P\rtimes\mathbb A)\times_{\mathbb B}(Y\rtimes\mathbb B)
\]
of the derived contextual execution object whenever the cospan computes the
homotopy pullback.

\item Up to the forward-opposite convention, these representatives recover the
strict relative program categories
\[
\mathsf{Prog}_{\Phi,Y}.
\]

\item A derived contextual program on
\[
\mathsf{Exec}^h_{\chi,\rho}
\]
is an endocorrespondence in
\[
\operatorname{Corr}(\mathcal C).
\]

\item Under a derived regular predicate theory, every correspondence
\[
R=(\mathbb X\xleftarrow{\ell}\mathbb R\xrightarrow{r}\mathbb Y)
\]
acts on predicates by
\[
\langle R\rangle\varphi
=
\Sigma_\ell(r^\ast\varphi).
\]

\item Under a derived first-order Heyting predicate theory, every such
correspondence also acts by
\[
[R]\varphi
=
\Pi_\ell(r^\ast\varphi).
\]

\item Tests satisfy
\[
\langle\theta?\rangle\varphi
=
\theta\wedge\varphi,
\]
and, when implication exists,
\[
[\theta?]\varphi
=
\theta\Rightarrow\varphi.
\]

\item Under suitable finite-join hypotheses, finite choice of subprograms
satisfies
\[
\langle U\vee V\rangle\varphi
=
\langle U\rangle\varphi
\vee
\langle V\rangle\varphi.
\]

\item Under suitable countable-coproduct hypotheses, every endoprogram
\[
U:\mathbb Z\nrightarrow\mathbb Z
\]
has a finite-path star
\[
U^\star_{\mathrm{path}}
=
\coprod_{n\ge0}U^n,
\]
satisfying
\[
U^\star_{\mathrm{path}}
\simeq
1_{\mathbb Z}\oplus U;U^\star_{\mathrm{path}}.
\]

\item This path star is the free monoid on \(U\) in
\[
\operatorname{Corr}(\mathcal C)(\mathbb Z,\mathbb Z)
\]
under the corresponding coproduct-distributivity hypotheses.

\item Under suitable continuity and completeness hypotheses,
\[
\langle U^\star_{\mathrm{path}}\rangle\varphi
=
\mu X.\bigl(\varphi\vee\langle U\rangle X\bigr),
\]
and, when boxes exist,
\[
[U^\star_{\mathrm{path}}]\varphi
=
\nu X.\bigl(\varphi\wedge[U]X\bigr).
\]

\item If a derived execution monoid
\[
\mathbf m:\mathbb Z\nrightarrow\mathbb Z
\]
is chosen, every subprogram
\[
U\to\mathbf m
\]
has a folded trace map
\[
U^\star_{\mathrm{path}}\to\mathbf m
\]
which folds finite paths into composite executions.
\end{enumerate}
\end{theorem}

\begin{proof}
Parts (1) and (2) are the definitions of derived Mealy morphisms and derived
downstream behaviours. Part (3) is
Theorem~\ref{thm:derived-contextual-restriction}. Parts (4) and (5) are
Definitions~\ref{def:derived-relative-execution-object} and
\ref{def:derived-input-downstream-projections}. Part (6) is
Proposition~\ref{prop:derived-terminal-trace-case}. Part (7) is
Proposition~\ref{prop:derived-stateless-contextual-restriction}. Parts (8)
and (9) are Theorem~\ref{thm:strict-relative-programs-as-derived-representatives}
and Corollary~\ref{cor:strict-restriction-derived-contextual}. Part (10) is
Definition~\ref{def:derived-contextual-program}. Parts (11) and (12) are
Definitions~\ref{def:derived-diamond-correspondence} and
\ref{def:derived-box-correspondence}. Part (13) follows from
Propositions~\ref{prop:derived-test-diamond-law} and
\ref{prop:derived-test-box-law}. Part (14) is
Proposition~\ref{prop:derived-finite-choice-law}. Part (15) is
Definition~\ref{def:derived-path-star} and
Proposition~\ref{prop:derived-path-star-equation}. Part (16) is
Theorem~\ref{thm:derived-path-star-free-monoid}. Part (17) is
Theorem~\ref{thm:derived-path-fixed-point-comparison}. Part (18) is
Definition~\ref{def:derived-folded-trace-map} and
Proposition~\ref{prop:meaning-derived-folded-trace}.
\end{proof}

\begin{remark}[Hierarchy of structure]
\label{rem:derived-relative-hierarchy}
The derived relative theory separates three layers.

First, the operational layer only uses finite limits in
\[
\mathcal C.
\]
A derived Mealy morphism and a downstream behaviour combine by homotopy
pullback:
\[
\mathbb X\times_{\mathbb B}\mathbb Y.
\]

Second, the program layer studies correspondences on this contextual execution
object:
\[
\mathbb Z\nleftarrow \mathbb U\nrightarrow \mathbb Z.
\]

Third, the logical layer requires a predicate theory. Diamonds require
dependent sums, boxes require dependent products, tests require finite
Heyting structure, finite choice requires finite joins, and path/fixed-point
comparisons require countable and continuity hypotheses.
\end{remark}

\begin{remark}[Final conceptual summary]
\label{rem:derived-relative-final-summary}
The strict relative theory says:
\[
\text{Mealy adapter}
+
\text{downstream action}
\quad\leadsto\quad
\text{coupled program category}.
\]
The derived relative theory says:
\[
\text{derived output cocycle}
+
\text{derived downstream behaviour}
\quad\leadsto\quad
\text{homotopy pullback execution object}.
\]
Dynamic proarrow logic is then the modal and fixed-point logic of
correspondences on this contextual execution object.

In the strict case, the primitive execution span is visible directly as the
arrow span of a program category. In the derived case, primitive programs are
additional endocorrespondences, induced by representatives or supplied as a
chosen derived execution monoid. This is the structural clarification provided
by the derived presentation.
\end{remark}